\subsection{Composition of systems of formal power series}

Let A be a commutative ring; we shall say that a system

\[
\mathbf {f} = (f _ {1}, \dots , f _ {p}) \in (\mathrm{A} [ [ \mathrm{X} _ {1}, \dots , \mathrm{X} _ {q} ] ]) ^ {p}\tag{8}
\]

of formal power series in the \(X_{j}\) ( \(1 \leqslant j \leqslant q\) ), with coefficients in A, is without constant term if this is true of all the \(f_{j}\) . For every system (8) of formal power series and every system

\[
\mathbf {g} = (g _ {1}, \dots , g _ {q}) \in (\mathrm{A} [ [ \mathrm{X} _ {1}, \dots , \mathrm{X} _ {r} ] ]) ^ {q}\tag{9}
\]

of \(q\) formal power series without constant term, we shall denote by \(\mathbf{f} \circ \mathbf{g}\) (or \(\mathbf{f}(\mathbf{g})\) ) the system of formal power series \(f_{j}(g_{1},\ldots ,g_{q})\) ( \(1\leqslant j\leqslant p\) ) in

\[
\left. \left(\mathrm{A} \left[ \left[ \mathrm{X} _ {1}, \dots , \mathrm{X} _ {r} \right] \right]\right) ^ {p} \right.
\]

(Algebra, Chapter IV, § 5, no. 5). If

\[
\mathbf {h} = (h _ {1}, \dots , h _ {r}) \in (\mathrm{A} [ [ \mathrm{X} _ {1}, \dots , \mathrm{X} _ {s} ] ]) ^ {r}
\]

is a third system without constant term, then

\[
(\mathbf {f} \circ \mathbf {g}) \circ \mathbf {h} = \mathbf {f} \circ (\mathbf {g} \circ \mathbf {h}).\tag{10}
\]

For every integer m,

\[
\left(\mathbf {f} ^ {(m)} \circ \mathbf {g} ^ {(m)}\right) \circ \mathbf {h} ^ {(m)} = \mathbf {f} ^ {(m)} \circ \left(\mathbf {g} ^ {(m)} \circ \mathbf {h} ^ {(m)}\right)
\]

where \(\mathbf{f}^{(m)}\) , \(\mathbf{g}^{(m)}\) , \(\mathbf{h}^{(m)}\) denote the systems of polynomials consisting of terms of total degree \(\leqslant m\) in the systems of formal power series f, g, h. But clearly the terms of total degree \(\leqslant m\) in the series of \((\mathbf{f} \circ \mathbf{g}) \circ \mathbf{h}\) (resp. \(\mathbf{f} \circ (\mathbf{g} \circ \mathbf{h})\) ) are the same as in \((\mathbf{f}^{(m)} \circ \mathbf{g}^{(m)}) \circ \mathbf{h}^{(m)}\) (resp. \(\mathbf{f}^{(m)} \circ (\mathbf{g}^{(m)} \circ \mathbf{h}^{(m)})\) ), whence our assertion.

For every system (8), we shall denote by \(M_{\mathbf{f}}\) or \(M_{\mathbf{f}}(\mathbf{X})\) , the Jacobian matrix \((\partial f_i / \partial \mathbf{X}_j)\) ( \(1 \leqslant i \leqslant p, 1 \leqslant j \leqslant q\) ) where \(i\) is the index of the rows and \(j\) that of the columns; for two systems (8) and (9), where \(\mathbf{g}\) is without constant term,

\[
M _ {\mathbf {f} \circ \mathbf {g}} = (M _ {\mathbf {f}} (\mathbf {g})). M _ {\mathbf {g}}\tag{11}
\]

where \(M_{\mathbf{f}}(\mathbf{g})\) is the matrix whose elements are obtained by substituting \(g_j\) for \(\mathbf{X}_j\) ( \(1 \leqslant j \leqslant q\) ) in each series element of \(M_{\mathbf{f}}\) ; this formula is just a reformulation of formula (9) of Algebra, Chapter IV, § 5, no. 8. We shall denote by \(M_{\mathbf{f}}(0)\) the matrix of constant terms of the elements of \(M_{\mathbf{f}}\) ; then we deduce from (11) that

\[
M _ {\mathbf {f} \circ \mathbf {g}} (0) = M _ {\mathbf {f}} (0). M _ {\mathbf {g}} (0).\tag{12}
\]

Given an integer \(n > 0\) , we shall write

\[
\mathbf {1} _ {n} = \mathbf {X} = (\mathrm{X} _ {1}, \dots , \mathrm{X} _ {n}) \in (\mathrm{A} [ [ \mathrm{X} _ {1}, \dots , \mathrm{X} _ {n} ] ]) ^ {n},\tag{13}
\]

which will be considered as a matrix with a single column.

For every system \(\mathbf{f} = (f_1, \ldots, f_n) \in (\mathrm{A}[[\mathrm{X}_1, \ldots, \mathrm{X}_n]])^n\) , \(M_{\mathbf{f}}\) is a square matrix of order \(n\) ; we shall denote by \(\mathrm{J}_{\mathbf{f}}\) or \(\mathrm{J}_{\mathbf{f}}(\mathbf{X})\) its determinant and by \(\mathrm{J}_{\mathbf{f}}(0)\) the constant term of \(\mathrm{J}_{\mathbf{f}}\) , equal to \(\det(M_{\mathbf{f}}(0))\) ; if \(\mathbf{g} = (g_1, \ldots, g_n)\) is a system without constant term in \((\mathrm{A}[[\mathrm{X}_1, \ldots, \mathrm{X}_n]])^n\) , then, by (11) and (12),

\[
J _ {f \circ g} = J _ {f} (g) \cdot J _ {g}\tag{14}
\]

\[
\mathrm{J} _ {\mathbf {f} \circ \mathbf {g}} (0) = \mathrm{J} _ {\mathbf {f}} (0) \mathrm{J} _ {\mathbf {g}} (0).\tag{15}
\]

PROPOSITION 5. Let A be a commutative ring and \(\mathbf{f} = (f_1, \ldots, f_n)\) a system without constant term of \(n\) series in \(\mathrm{A}[[\mathrm{X}_1, \ldots, \mathrm{X}_n]]\) . Suppose that \(\mathrm{J}_{\mathbf{f}}(0)\) is invertible in A. Then there exists a system without constant term \(\mathbf{g} = (g_1, \ldots, g_n)\) of \(n\) series in \(\mathrm{A}[[\mathrm{X}_1, \ldots, \mathrm{X}_n]]\) such that

\[
\mathbf {f} \circ \mathbf {g} = \mathbf {1} _ {n}.\tag{16}
\]

This system is unique and

\[
\mathbf {g} \circ \mathbf {f} = \mathbf {1} _ {n}.\tag{17}
\]

The existence and uniqueness of \(\mathbf{g}\) follows from Algebra, Chapter IV, § 5, no. 9, Proposition 10, applied to the \(n\) formal power series

\[
f _ {i} (\mathrm{Y} _ {1}, \dots , \mathrm{Y} _ {n}) - \mathrm{X} _ {i} \quad (1 \leqslant i \leqslant n).
\]

It follows from (15) and (16) that \(\mathrm{J}_{\mathbf{f}}(0)\mathrm{J}_{\mathbf{g}}(0)=1\) and hence \(\mathrm{J}_{\mathbf{g}}(0)\) is also invertible. We conclude that there exists a system \(\mathbf{h}=(h_{1},\ldots,h_{n})\) of n series without constant term in \(\mathrm{A}[[\mathrm{X}_{1},\ldots,\mathrm{X}_{n}]]\) such that \(g\circ h=1_{n}\) ; from this relation and (16) it then follows, with the aid of (10), that

\[
\mathbf {h} = \mathbf {1} _ {n} \circ \mathbf {h} = (\mathbf {f} \circ \mathbf {g}) \circ \mathbf {h} = \mathbf {f} \circ (\mathbf {g} \circ \mathbf {h}) = \mathbf {f} \circ \mathbf {1} _ {n} = \mathbf {f}.
\]

Proposition 5 and formulae (10) and (15) show that under the law of composition \((\mathbf{f},\mathbf{g})\mapsto \mathbf{f}\circ \mathbf{g}\) the set of systems \(\mathbf{f} = (f_1,\dots ,f_n)\) of \(n\) series without constant term in \(\mathrm{A}[[\mathrm{X}_1,\dots ,\mathrm{X}_n]]\) for which \(\mathrm{J}_{\mathbf{f}}(0)\) is invertible in A, is a group.

\subsection{Systems of equations in complete rings}

To abbreviate, we shall say in what follows that a ring satisfies Hensel's conditions if it is commutative, linearly topologized, Hausdorff and complete; given an ideal m in such a ring, m (or the ordered pair (A, m)) will be said to satisfy Hensel's conditions if m is closed in A and its elements are topologically nilpotent. The ideal t of A consisting of all the topologically nilpotent elements satisfies Hensel's conditions (no. 3).

In particular, if A is a commutative ring and m an ideal of A and A is Hausdorff and complete with respect to the m-adic topology, the ordered pair (A, m) satisfies Hensel's conditions.

PROPOSITION 6. Let A be a commutative ring, B a ring satisfying Hensel's conditions and \(u: \mathbf{A} \to \mathbf{B}\) a homomorphism. For every family \(\mathbf{x} = (x_1, \ldots, x_n)\) of topologically nilpotent elements of B, there exists a unique homomorphism \(\tilde{u}\) from \(\mathbf{A}[[\mathbf{X}_1, \ldots, \mathbf{X}_n]]\) to B such that \(\tilde{u}(a) = u(a)\) for all \(a \in \mathbf{A}\) and \(\tilde{u}(\mathbf{X}_i) = x_i\) for \(1 \leqslant i \leqslant n\) . Moreover, if m denotes the ideal of series without constant term in \(\mathbf{A}[[\mathbf{X}_1, \ldots, \mathbf{X}_n]]\) , \(\tilde{u}\) is continuous for the m-adic topology.

Let \(\mathfrak{a}\) be the finitely generated ideal generated in \(\mathbf{B}\) by the \(x_{i}\) ( \(1 \leqslant i \leqslant n\) ); for every open ideal \(\mathfrak{S}\) of \(\mathbf{B}\) , the images of the \(x_{i}\) in \(\mathbf{B}/\mathfrak{S}\) are nilpotent, hence the ideal \((\mathfrak{a} + \mathfrak{S})/\mathfrak{S}\) is nilpotent in \(\mathbf{B}/\mathfrak{S}\) and there exists an integer \(k\) such that, for \(\sum_{i=1}^{n} p_{i} \geqslant k\) , \(x_{1}^{p_{1}} \ldots x_{n}^{p_{n}} \in \mathfrak{S}\) . As every element of \(m^{k}\) is a finite sum of formal power series of the form \(X_{1}^{p_{1}} \ldots X_{n}^{p_{n}} g(X_{1}, \ldots, X_{n})\) , where \(\sum_{i=1}^{n} p_{i} \geqslant k\) , it is seen that, if \(\tilde{u}\) solves the problem, then \(\tilde{u}(\mathfrak{m}^k) \subset \mathfrak{S}\) , which proves the continuity of \(\tilde{u}\) . There obviously exists a unique homomorphism

\[
v \colon \mathrm{A} [ \mathrm{X} _ {1}, \dots , \mathrm{X} _ {n} ] \rightarrow \mathrm{B}
\]

such that \(v(a) = u(a)\) for \(a \in \mathbf{A}\) and \(v(\mathbf{X}_i) = x_i\) for \(1 \leqslant i \leqslant n\) and the above argument shows that \(v\) is continuous with respect to the topology induced on \(\mathbf{A}[\mathbf{X}_1, \ldots, \mathbf{X}_n]\) by the m-adic topology. As \(\mathbf{A}[\mathbf{X}_1, \ldots, \mathbf{X}_n]\) is dense in \(\mathbf{A}[[\mathbf{X}_1, \ldots, \mathbf{X}_n]]\) with the m-adic topology and B is Hausdorff and complete, this completes the proof of the existence and uniqueness of \(\tilde{u}\) .

Note that this proposition gives us again as a special case (i) of Proposition 11 of § 2, no. 9.

If A itself is linearly topologized, the restriction of \(\tilde{u}\) to \(A\{X_{1},\ldots,X_{n}\}\) coincides with the homomorphism derived from u in Proposition 4 of no. 2. This follows immediately from the fact that \(A[X_{1},\ldots,X_{n}]\) is dense in \(A\{X_{1},\ldots,X_{n}\}\) if this ring is given the topology with fundamental system of neighbourhoods of 0 the ideals \(m^{k}\cap N_{3}\) (in the notation of no. 2, this topology is the least upper bound of the topology induced on \(A\{X_{1},\ldots,X_{n}\}\) by the m-adic topology on \(A[[X_{1},\ldots,X_{n}]]\) and the topology defined in no. 2).

If B = A and u is the identity mapping, we shall write \(f(x_{1}, \ldots, x_{n})\) or \(f(\mathbf{x})\) for the element \(\tilde{u}(f)\) for every formal power series \(f \in A[[X_{1}, \ldots, X_{n}]]\) ; for every system \(\mathbf{f} = (f_{1}, \ldots, f_{r})\) of formal power series of \(A[[X_{1}, \ldots, X_{n}]]\) , let \(\mathbf{f}(\mathbf{x})\) denote the element \((f_{1}(\mathbf{x}), \ldots, f_{r}(\mathbf{x}))\) of \(A^{r}\) , then it is said to be obtained by substituting the \(x_{i}\) for the \(X_{i}\) in f.~If \(n \leqslant m\) and F is a formal power series of \(A[[X_{1}, \ldots, X_{m}]]\) , it is possible to consider F as a formal power series in \(X_{n+1}, \ldots, X_{m}\) with coefficients in \(A[[X_{1}, \ldots, X_{n}]]\) ; let

\[
\mathbf {F} (x _ {1}, \dots , x _ {n}, \mathbf {X} _ {n + 1}, \dots , \mathbf {X} _ {m})
\]

denote the formal power series in \(\mathrm{A}[[\mathrm{X}_{n+1}, \ldots, \mathrm{X}_{m}]]\) obtained by substituting the x \(_{i}\) for the X \(_{i}\) in the coefficients of F, for 1 ≤ i ≤ n.

Let us take B to be the ring of formal power series \(\Lambda[[\mathbf{X}_1,\ldots,\mathbf{X}_r]]\) and let \(n\) be the ideal of series in B without constant term, so that (B, n) satisfies Hensel's conditions (§ 2, no. 6, Corollary to Proposition 6). Proposition 6 may be applied by taking the \(x_i \in B\) to be series without constant term; then, for every series \(f \in A[[\mathbf{X}_1,\ldots,\mathbf{X}_n]]\) , \(\tilde{u}(f)\) is just the formal power series \(f(x_1,\ldots,x_n)\) defined in Algebra, Chapter IV, § 5, no. 5. This is obvious if \(f\) is a polynomial and it follows from the proposition in the general case by remarking that \(f \mapsto f(x_1,\ldots,x_n)\) is continuous on \(A[[\mathbf{X}_1,\ldots,\mathbf{X}_n]]\) with respect to the m-adic topology.

COROLLARY. Let A be a ring satisfying Hensel's condition and \(\mathbf{x} = (x_{1},\ldots ,x_{n})\) a family of topologically nilpotent elements of A. Let \(\mathbf{g} = (g_1,\dots,g_q)\) be a system without constant term of series in \(\mathbf{A}[[\mathbf{X}_1,\ldots ,\mathbf{X}_n]]\) and \(\mathbf{f} = (f_1,\dots,f_p)\) a system of formal power series in \(\mathbf{A}[[\mathbf{X}_1,\dots,\mathbf{X}_q]]\) . Then \(\mathbf{g}(\mathbf{x}) = (g_1(\mathbf{x}),\dots,g_q(\mathbf{x}))\) is a family of topologically nilpotent elements of \(\mathbf{A}\) and

\[
(\mathbf {f} \circ \mathbf {g}) (\mathbf {x}) = \mathbf {f} (\mathbf {g} (\mathbf {x})).\tag{18}
\]

The fact that the \(g_{i}(\mathbf{x})\) are topologically nilpotent follows immediately from Proposition 6 and the fact that in A the ideal of topologically nilpotent elements is closed. Relation (18) is obvious when the \(f_{j}\) are polynomials; on the other hand, if m and \(m'\) are the ideals of series without constant term in

\[
\mathrm{A} \left[ \left[ \mathrm{X} _ {1}, \dots , \mathrm{X} _ {q} \right] \right] \quad \text { and } \quad \mathrm{A} \left[ \left[ \mathrm{X} _ {1}, \dots , \mathrm{X} _ {n} \right] \right]
\]

respectively, clearly the relation \(f \in \mathfrak{m}^k\) implies \(f(g_1, \ldots, g_q) \in \mathfrak{m}'^k\) . The two sides of (18) are therefore continuous functions of \(\mathbf{f}\) to \((\mathrm{A}[[\mathrm{X}_1, \ldots, \mathrm{X}_q]])^p\) if \(\mathrm{A}[[\mathrm{X}_1, \ldots, \mathrm{X}_q]]\) is given the m-adic topology, by virtue of the above remark and Proposition 6; whence relation (18).

In what follows, for a ring A and an ideal m of A we shall denote by \(m^{\times n}\) the

product set \(\prod_{i=1}^{n} m_i\) in \(A^n\) , where \(m_i = m\) for \(1 \leqslant i \leqslant n\) , to avoid ambiguity.

PROPOSITION 7. Let A be a ring and m an ideal of A such that the ordered pair (A, m) satisfies Hensel's conditions. Let \(\mathbf{f} = (f_1, \ldots, f_n)\) be a system without constant term of series in \(\mathrm{A}[[\mathrm{X}_1, \ldots, \mathrm{X}_n]]\) such that \(\mathrm{J}_{\mathbf{f}}(0)\) is invertible in A. Then, for all \(\mathbf{x} \in \mathfrak{m}^{\times n}\) , \(\mathbf{f}(\mathbf{x}) \in \mathfrak{m}^{\times n}\) and \(\mathbf{x} \mapsto \mathbf{f}(\mathbf{x})\) is a bijection of \(\mathfrak{m}^{\times n}\) onto itself, the inverse bijection being \(\mathbf{x} \mapsto \mathbf{g}(\mathbf{x})\) , where \(\mathbf{g}\) is given by relation (16) in no. 4.

The fact that \(\mathbf{f}(\mathbf{x})\in\mathfrak{m}^{\times n}\) is obvious when the \(f_{i}\) are polynomials and follows in the general case from Proposition 6 and the fact that m is closed in A. The other assertions of the proposition are then immediate consequences of (16), (17) and (18).

COROLLARY. Let \(\mathfrak{q}\) be a closed ideal of A contained in \(\mathfrak{m}\) . Then the relation \(\mathbf{x} \equiv \mathbf{x}'\) (mod. \(\mathfrak{q}^{\times n}\) ) is equivalent to \(\mathbf{f}(\mathbf{x}) \equiv \mathbf{f}(\mathbf{x}')\) (mod. \(\mathfrak{q}^{\times n}\) ).

For every formal power series \(f \in \mathbf{A}[[\mathbf{X}_1, \ldots, \mathbf{X}_n]]\) ,

\[
f \left(\mathrm{X} _ {1}, \dots , \mathrm{X} _ {n}\right) - f \left(\mathrm{Y} _ {1}, \dots , \mathrm{Y} _ {n}\right) = \sum_ {i = 1} ^ {n} \left(\mathrm{X} _ {i} - \mathrm{Y} _ {i}\right) h _ {i} \left(\mathrm{X} _ {1}, \dots , \mathrm{X} _ {n}, \mathrm{Y} _ {1}, \dots , \mathrm{Y} _ {n}\right)
\]

where the \(h_i\) belong to \(\mathbf{A}[[\mathbf{X}_1,\ldots,\mathbf{X}_n,\mathbf{Y}_1,\ldots,\mathbf{Y}_n]]\) (Algebra, Chapter IV, § 5, no. 8, Proposition 9); it follows immediately that the relation \(\mathbf{x} \equiv \mathbf{x}'\) (mod. \(\mathfrak{q}^{\times n}\) ) implies \(\mathbf{f}(\mathbf{x}) \equiv \mathbf{f}(\mathbf{x}')\) (mod. \(\mathfrak{q}^{\times n}\) ). The converse is obtained by replacing \(\mathbf{f}\) by its ``inverse'' \(\mathbf{g}\) .

THEOREM 2. Let \(\mathbf{A}\) be a ring and \(\mathfrak{m}\) an ideal of \(\mathbf{A}\) such that the ordered pair \((\mathbf{A},\mathfrak{m})\) satisfies Hensel's conditions. Let \(\mathbf{f} = (f_1,\ldots ,f_n)\) be a system of \(n\) elements of \(\mathbf{A}\{\mathbf{X}_1,\ldots ,\mathbf{X}_n\}\) and let \(\mathbf{a}\in \mathbf{A}^n\) ; let us write \(\mathrm{J}_{\mathbf{f}}(\mathbf{a}) = e\) . There exists a system \(\mathbf{g} = (g_1,\dots ,g_n)\) of restricted formal power series without constant term in \(\mathbf{A}\{\mathbf{X}_1,\dots ,\mathbf{X}_n\}\) such that

\begin{enumerate}
\def\labelenumi{(\roman{enumi})}
\item
  \(M_{\mathbf{g}}(0) = I_{n}\) (unit matrix).
\item
  For all \(\mathbf{x} \in \mathbf{A}^n\) ,

\[
\mathbf {f} (\mathbf {a} + e \mathbf {x}) = \mathbf {f} (\mathbf {a}) + M _ {\mathbf {f}} (\mathbf {a}). e \mathbf {g} (\mathbf {x}).\tag{19}
\]

\item
  Let \(\mathbf{h} = (h_1, \ldots, h_n)\) be the system of formal power series without constant term (not necessarily restricted) such that \(\mathbf{g} \circ \mathbf{h} = \mathbf{1}_n\) (Proposition 5). For all \(y \in \mathfrak{m}^{\times n}\) ,

\[
\mathbf {f} (\mathbf {a} + e \mathbf {h} (\mathbf {y})) = \mathbf {f} (\mathbf {a}) + M _ {\mathbf {f}} (\mathbf {a}). e \mathbf {y}.\tag{20}
\]

\end{enumerate}

For every formal power series \(f \in \mathbf{A}[[\mathbf{X}_1, \ldots, \mathbf{X}_n]]\) ,

\[
f (\mathbf {X} + \mathbf {Y}) = f (\mathbf {X}) + M _ {f} (\mathbf {X}). \mathbf {Y} + \sum_ {1 \leqslant i \leqslant j \leqslant n} \mathrm{G} _ {i j} (\mathbf {X}, \mathbf {Y}) \mathrm{Y} _ {i} \mathrm{Y} _ {j}\tag{21}
\]

where the \(G_{ij}\) are well determined formal power series in

\[
\mathrm{A} \left[ \left[ \mathrm{X} _ {1}, \dots , \mathrm{X} _ {n}, \mathrm{Y} _ {1}, \dots , \mathrm{Y} _ {n} \right] \right].
\]

If \(f\) is restricted, so are the elements of \(M_{f}\) and the \(G_{ij}\) , for these formal power series are polynomials if \(f\) is a polynomial and it follows from their uniqueness that for every open ideal \(\Im\) of \(A\) , denoting by \(p_{\Im} \colon A \to A / \Im\) the canonical mapping, the image of \(G_{ij}\) under \(\bar{p}_{\Im}\) is the coefficient of \(Y_{i}Y_{j}\) in \(\bar{p}_{\Im}(F)\) where \(F\) is the formal power series \(f(\mathbf{X} + \mathbf{Y})\) in \(A[[X_1, \ldots, X_n, Y_1, \ldots, Y_n]]\) ; whence our assertion.

This being so, writing formula (21) for each series \(f_{i}\) ( \(1 \leqslant i \leqslant n\) ), we obtain for all \(\mathbf{x} \in \mathbf{A}^n\) (no. 2, Proposition 4),

\[
\mathbf {f} (\mathbf {a} + e \mathbf {x}) = \mathbf {f} (\mathbf {a}) + M _ {\mathbf {f}} (\mathbf {a}). e \mathbf {x} + e ^ {2} \mathbf {r} (\mathbf {x})\tag{22}
\]

where \(\mathbf{r} = (r_1, \ldots, r_n)\) is a system of restricted formal power series each of which is of total order \(\geqslant 2\) . It follows from formulae (18) of Algebra, Chapter III, §6, no. 5 that there exists a square matrix \(M' \in \mathbf{M}_n(\mathbf{A})\) such that

\[
M _ {\mathbf {f}} (\mathbf {a}). M ^ {\prime} = e I _ {n},\tag{23}
\]

whence using this in (22)

\[
\mathbf {f} (\mathbf {a} + e \mathbf {x}) = \mathbf {f} (\mathbf {a}) + M _ {\mathbf {f}} (\mathbf {a}). e \mathbf {x} + M _ {\mathbf {f}} (\mathbf {a}) M ^ {\prime}. e \mathbf {r} (\mathbf {x}).\tag{24}
\]

Writing \(\mathbf{g} = \mathbf{1}_n + M'.\mathbf{r}\) , we see that \(\mathbf{g}\) satisfies conditions (i) and (ii); then it is sufficient to replace \(\mathbf{x}\) by \(\mathbf{h}(\mathbf{y})\) to obtain (iii).

COROLLARY 1. Let A be a ring and m an ideal of A such that the ordered pair (A, m) satisfies Hensel's conditions. Let \(f \in \mathrm{A}\{\mathbf{X}\}\) , \(a \in \mathrm{A}\) and write \(e = f'(a)\) . If \(f(a) \equiv 0 \pmod{e^2\mathfrak{m}}\) , then there exists \(b \in \mathrm{A}\) such that \(f(b) = 0\) and \(b \equiv a \pmod{e\mathfrak{m}}\) . If \(b'\) is another element of A such that \(f(b') = 0\) and \(b' \equiv a (\text{mod. em})\) , then \(e(b - b') = 0\) . In particular, b is unique if e is not a divisor of zero in A.

Let \(f(a) = e^{2}c\) where \(c \in \mathfrak{m}\) ; formula (20) for \(n = 1\) gives

\[
f (a + e h (y)) = e ^ {2} (c + y)
\]

and it is therefore sufficient to take \(y = -c\) , \(b = a + eh(-c)\) . Moreover if \(b = a + ex\) , \(b' = a + ex'\) , \(x \in \mathfrak{m}\) , \(x' \in \mathfrak{m}\) , \(f(b) = f(b') = 0\) , we deduce from (19) that \(e^2(g(x) - g(x')) = 0\) . As \(g(\mathbf{X}) - g(\mathbf{Y}) = (\mathbf{X} - \mathbf{Y})u(\mathbf{X}, \mathbf{Y})\) , where \(u\) is restricted and \(u(0,0) = 1\) , \(g(x) - g(x') = (x - x')v\) , where \(v \in \mathbf{A}\) is invertible, for, since \(\mathfrak{m}\) is closed, \(v - 1 = u(x,x') - 1 \in \mathfrak{m}\) and \(\mathfrak{m}\) is contained in the Jacobson radical of \(\mathbf{A}\) ; whence the relation \(e(b - b') = 0\) .

Remark. The corollary applies notably when \(e\) is invertible in A; we can then also deduce the existence of \(b\) from Hensel's Theorem, for the canonical image of \(f(\mathbf{X})\) in \((\mathbf{A} / \mathfrak{m})\{\mathbf{X}\}\) is of the form \((\mathbf{X} - \alpha)f_1(\mathbf{X}), \mathbf{X} - \alpha\) and \(f_1(\mathbf{X})\) being strongly relatively prime, for \(f_1(\alpha) = f'(\alpha)\) is the image of \(e\) (no. 1, Example).

\textbf{Examples}\par

\begin{enumerate}
\def\labelenumi{(\arabic{enumi})}
\item
  Let \(p\) be a prime number \(\neq 2\) and \(n\) an integer whose class mod. \(p\) is a square \(\neq 0\) in the prime field \(\mathbf{F}_p\) . If \(\mathbf{Z}_p\) is the ring of \(p\) -adic integers (§ 2, no. 12, Example 3), the application of Corollary 1 to the polynomial \(\mathbf{X}^2 - n\) shows that \(n\) is a square in \(\mathbf{Z}_p\) ; for example 7 is a square in \(\mathbf{Z}_3\) .
\item
  Let \(\mathbf{A} = \mathbf{K}[[\mathbf{Y}]]\) be the ring of formal power series in one indeterminate with coefficients in a commutative field \(\mathbf{K}\) ; with the (Y)-adic topology, the ring \(\mathbf{A}\) is Hausdorff and complete (§ 2, no. 6, Corollary to Proposition 6) and the mapping \(f(\mathbf{Y}) \mapsto f(0)\) defines by passing to the quotient ring an isomorphism of \(\mathbf{A}/(\mathbf{Y})\) onto the field \(\mathbf{K}\) . By Corollary 1, if \(\mathbf{F}(\mathbf{Y}, \mathbf{X})\) is a polynomial in \(\mathbf{X}\) with coefficients in \(\mathbf{A}\) and \(a\) is a simple root of \(\mathbf{F}(0, \mathbf{X})\) in \(\mathbf{K}\) , there exists a unique formal power series \(f(\mathbf{Y})\) such that \(f(0) = a\) and \(\mathbf{F}(\mathbf{Y}, f(\mathbf{Y})) = 0\) .
\end{enumerate}

COROLLARY 2. Let A be a ring and m an ideal of A such that the ordered pair (A, m) satisfies Hensel's conditions. Let r, n be integers such that \(0 \leqslant r < n\) and \(\mathbf{f} = (f_{r+1}, \ldots, f_n)\) is a system of n - r elements of \(A\{X_1, \ldots, X_n\}\) ; let \(\mathrm{J}_{\mathbf{f}}^{(n-r)}(\mathbf{X})\) denote the minor of \(M_{\mathbf{f}}(\mathbf{X})\) consisting of the columns of index j such that \(r + 1 \leqslant j \leqslant n\) . Let \(a \in A^n\) be such that \(\mathrm{J}_{\mathbf{f}}^{(n-r)}(\mathbf{a})\) is invertible in A and \(\mathbf{f}(\mathbf{a}) \equiv 0 \pmod{\mathfrak{m}^{\times(n-r)}}\) . Then there exists a unique \(\mathbf{x} = (x_1, \ldots, x_n) \in \mathrm{A}^n\) such that \(x_k = a_k\) for \(1 \leqslant k \leqslant r\) , \(x \equiv a \pmod{m^{\times n}}\) and \(\mathbf{f}(\mathbf{x}) = 0\) .

Substituting \(a_k\) for \(\mathbf{X}_k\) for \(1 \leqslant k \leqslant r\) in the \(f_i\) (no. 2, Remark 3), we see immediately that we may restrict our attention to the case where \(r = 0\) to prove the corollary. Theorem 1 and Proposition 7 show then that \(\mathbf{f}\) defines a bijection of \(\mathbf{a} + \mathfrak{m}^{\times n}\) onto \(\mathbf{f}(\mathbf{a}) + \mathfrak{m}^{\times n} = \mathfrak{m}^{\times n}\) ; the corollary follows from the fact that \(0 \in \mathfrak{m}^{\times n}\) .

COROLLARY 3. With the notation of Corollary 2, let \(\mathbf{a} \in \mathbf{A}^n\) ; let us write \(e = \mathrm{J}_{\mathbf{f}}^{(n-r)}(\mathbf{a})\) (not necessarily invertible in A) and suppose that \(\mathbf{f}(\mathbf{a}) \equiv 0 \pmod{e^2\mathfrak{m}^{\times(n-r)}}\) . Then there exist \(n - r\) formal power series without constant term \(\phi_i (r + 1 \leqslant i \leqslant n)\) in \(\mathbf{A}[[\mathbf{X}_1, \ldots, \mathbf{X}_r]]\) such that, for all \(\mathbf{t} = (t_1, \ldots, t_r) \in \mathfrak{m}^{\times r}\) ,

\[
f _ {i} \left(a _ {1} + e ^ {2} t _ {1}, \dots , a _ {r} + e ^ {2} t _ {r}, a _ {r + 1} + e \phi_ {r + 1} (\mathbf {t}), \dots , a _ {n} + e \phi_ {n} (\mathbf {t})\right) = 0
\]

\[
\text{for } r + 1 \leqslant i \leqslant n.
\]

For \(1 \leqslant i \leqslant r\) , let \(f_{i}(\mathbf{X}) = \mathbf{X}_{i} - a_{i}\) and let \(\mathbf{u} = (f_{1}, \ldots, f_{n})\) ; then \(\mathrm{J}_{\mathbf{u}}(\mathbf{a}) = e\) and Theorem 2 may be applied to the system \(\mathbf{u}\) . With the notation of Theorem 2 it follows from the above definitions that \(g_{i}(\mathbf{X}) = \mathbf{X}_{i}\) for \(1 \leqslant i \leqslant r\) , whence \(h_{i}(\mathbf{X}) = \mathbf{X}_{i}\) for \(1 \leqslant i \leqslant r\) ; moreover, if \(M' \in \mathbf{M}_{n}(\mathbf{A})\) is such that \(M_{\mathbf{u}}(\mathbf{a}) \cdot M' = eI_{n}\) , \(M'\) is of the form

\[
\left( \begin{array}{c c} e I _ {r} & 0 \\ * & * \end{array} \right).
\]

Replacing \(\mathbf{y}\) by \(M^{\prime}.\mathbf{z}\) (where \(\mathbf{z} = (z_{1},\ldots ,z_{n})\in \mathfrak{m}^{\times n}\) ) in formula (20), we obtain

\[
\begin{array}{r l} f _ {i} (a _ {1} + e ^ {2} z _ {1}, \dots , a _ {r} + e ^ {2} z _ {r}, a _ {r + 1} + e h _ {r + 1} (M ^ {\prime}. \mathbf {z}), \dots , a _ {n} + e h _ {n} (M ^ {\prime}. \mathbf {z})) \\ & = f _ {i} (a) + e ^ {2} z _ {i} \quad \text { for } \quad 1 \leqslant i \leqslant n. \end{array}\tag{26}
\]

By hypothesis, \(f_{j}(a) = e^{2}b_{j}\) where \(b_{j} \in \mathfrak{m}\) for \(r + 1 \leqslant j \leqslant n\) . Let us write \(\psi_{j}(\mathbf{X}_{1},\ldots ,\mathbf{X}_{n}) = h_{j}(M^{\prime}.\mathbf{X})\) and

\[
\phi_ {j} \left(\mathrm{X} _ {1}, \dots , \mathrm{X} _ {r}\right) = \psi_ {j} \left(\mathrm{X} _ {1}, \dots , \mathrm{X} _ {r}, - b _ {r + 1}, \dots , - b _ {n}\right)
\]

for \(r + 1 \leqslant j \leqslant n\) . For \(r + 1 \leqslant i \leqslant n\) , substituting \(t_j\) for \(z_j\) for \(1 \leqslant j \leqslant r\) and \(b_j\) for \(z_j\) for \(r + 1 \leqslant j \leqslant n\) in (26), we obtain relations (25) for all \(t \in m^{\times r}\) .

\subsection{Application to decompositions of rings}

LEMMA 2. Let A be a ring and m an ideal of A such that the ordered pair (A, m) satisfies Hensel's conditions. Let B be the quotient ring A/m and \(\pi: A \to B\) be the canonical homomorphism. For every idempotent c of B there exists a unique idempotent e of A such that \(\pi(e) = c\) .

Let \(a\) be such that \(\pi(a) = c\) ; Corollary 1 to Theorem 2 of no. 5 may be applied to the polynomial \(f(\mathbf{X}) = \mathbf{X}^2 - \mathbf{X}\) in \(\mathbf{A}[\mathbf{X}]\) and the element \(a \in \mathbf{A}\) . Then \(f'(a) = 2a - 1\) and, as \(\pi(2a - 1) = 2c - 1\) and \((2c - 1)^2 = 1\) in \(\mathbf{B}\) , \(2c - 1\) is invertible in \(\mathbf{B}\) and hence \(2a - 1\) is invertible in \(\mathbf{A}\) (§ 2, no. 13, Lemma 3). As \(f(a) \in \mathfrak{m}\) , Corollary 1 to Theorem 2 of no. 5 immediately gives the existence and uniqueness of \(e\) .

PROPOSITION 8. Let A be a ring and m an ideal of A such that the ordered pair (A, m) satisfies Hensel's conditions. Let B be the quotient ring A/m and π: A → B the canonical homomorphism. If B is the direct composition of a finite family \((\mathfrak{b}_{i})_{i\in I}\) of ideals, there exists a unique family \((\mathfrak{a}_{i})_{i\in I}\) of ideals of A such that \(\pi(\mathfrak{a}_{i}) = \mathfrak{b}_{i}\) for all \(i \in I\) and A is the direct composition of the family \((\mathfrak{a}_{i})\) .

Let \(1 = \sum_{i} c_{i}\) where \(c_{i} \in \mathfrak{b}_{i}\) for all \(i\) ; the \(c_{i}\) are idempotents of \(\mathbf{B}\) such that \(c_{i}c_{j} = 0\) for \(i \neq j\) . By Lemma 2 there therefore exist idempotents \(e_{i}\) of \(\mathbf{A} (i \in \mathbf{I})\) such that \(\pi(e_{i}) = c_{i}\) for all \(i\) ; as \(e_{i}e_{j}\) is an idempotent such that

\[
\pi (e _ {i} e _ {j}) = c _ {i} c _ {j} = 0
\]

for \(i \neq j\) , \(e_i e_j = 0\) for \(i \neq j\) (Lemma 2); as \(1 - \sum_{i} e_i\) is an idempotent such that \(\pi\left(1 - \sum_{i} e_i\right) = 1 - \sum_{i} c_i = 0\) , similarly \(1 = \sum_{i} e_i\) . It follows that \(A\) is the direct composition of the ideals \(a_i = e_i A\) and that \(\pi(a_i) = \pi(e_i) B = b_i\) .

It remains to show the uniqueness of such a decomposition. Now, suppose that A is the direct composition of another family \((\mathfrak{a}_i')_{i\in I}\) of ideals such that \(\pi (\mathfrak{a}_i^{\prime}) = \mathfrak{b}_i\) for all \(i\) ; then \(1 = \sum_{i}e_{i}^{\prime}\) where \(e_i^\prime \in \mathfrak{a}_i^\prime\) , whence, in B, \(1 = \sum_{i}\pi (e_i^{\prime})\) where \(\pi (e_i^{\prime})\in \mathfrak{b}_i\) , which implies \(\pi (e_i^{\prime}) = c_i\) ; as \(e_i^\prime\) and \(e_i\) are idempotents, necessarily \(e_i^\prime = e_i\) (Lemma 2), which completes the proof.

Remark. Proposition 8 again gives the structure of a semi-local ring on A which is Hausdorff and complete with the r-adic topology (r the Jacobson radical of A), which has already been obtained as a consequence of § 2, no. 13, Corollary to Proposition 19.

\section{FLATNESS PROPERTIES OF FILTERED MODULES}

\subsection{Ideally Hausdorff modules}

DEFINITION 1. Let A be a commutative ring and \(\mathfrak{I}\) an ideal of A. An A-module M is called ideally Hausdorff with respect to \(\mathfrak{I}\) (or simply ideally Hausdorff if there is no ambiguity) if, for every finitely generated ideal \(\mathfrak{a}\) of A, the A-module \(\mathfrak{a} \otimes_{\mathbf{A}} \mathbf{M}\) is Hausdorff with the \(\mathfrak{I}\) -adic topology.

Putting \(\mathfrak{a} = \mathbf{A}\) in this definition, we have already seen that \(\mathbf{M}\) is necessarily Hausdorff with the 3-adic topology.

\textbf{Examples}\par

\begin{enumerate}
\def\labelenumi{(\arabic{enumi})}
\item
  If A is Noetherian and \(\mathfrak{I}\) is contained in the Jacobson radical of A (in other words if A is a Zariski ring with the \(\mathfrak{I}\) -adic topology), every finitely generated A-module is ideally Hausdorff ( \(\S\) 3, no. 3, Proposition 6).
\item
  Every direct sum of ideally Hausdorff modules is an ideally Hausdorff module, by virtue of the relations
\end{enumerate}

\[
\Im^ {n} \left(a \otimes_ {A} \bigoplus_ {\lambda \in L} M _ {\lambda}\right) = \Im^ {n} \bigoplus_ {\lambda \in L} (a \otimes_ {A} M _ {\lambda}) = \bigoplus_ {\lambda \in L} \Im^ {n} (a \otimes_ {A} M _ {\lambda}).
\]

\begin{enumerate}
\def\labelenumi{(\arabic{enumi})}
\setcounter{enumi}{2}
\tightlist
\item
  If an A-module M is flat and Hausdorff with the \(\Im\) -adic topology it is ideally Hausdorff, for \(a \otimes_{A} M\) is then identified with a submodule of M and the \(\Im\) -adic topology on \(a \otimes_{A} M\) is finer than the topology induced on \(a \otimes_{A} M\) by the \(\Im\) -adic topology on M, which is Hausdorff by hypothesis.
\end{enumerate}

\subsection{Statement of the flatness criterion}

Let A be a ring, \(\mathfrak{I}\) a two-sided ideal of A, M a left-module and gr(A) and gr(M) the graded ring and graded gr(A)-module associated respectively with the ring A and with the module M with the \(\mathfrak{I}\) -adic filtrations (§ 2, no. 3). We have seen (loc. cit.) that for every integer \(n \geqslant 0\) there is a surjective Z-module homomorphism

\[
\gamma_ {n} \colon (\mathfrak {I} ^ {n} / \mathfrak {I} ^ {n + 1}) \otimes_ {\mathrm{A} / \mathfrak {I}} (\mathrm{M} / \mathfrak {I} \mathrm{M}) \rightarrow \mathfrak {I} ^ {n} \mathrm{M} / \mathfrak {I} ^ {n + 1} \mathrm{M}
\]

and a graded homomorphism of degree 0 of graded gr(A)-modules

\[
\gamma_ {M} \colon \operatorname{gr} (A) \otimes_ {\mathbf {g r} _ {0} (A)} \operatorname{gr} _ {0} (M) \to \operatorname{gr} (M)
\]

whose restriction to \(\mathrm{gr}_n(\mathbf{A})\otimes_{\mathbf{gr}_0(\mathbf{A})}\mathrm{gr}_0(\mathbf{M})\) is \(\gamma_{n}\) for all \(n\) and which is therefore surjective.

THEOREM 1. Let A be a commutative ring, \(\Im\) an ideal of A and M an A-module. Consider the following properties:

\begin{enumerate}
\def\labelenumi{(\roman{enumi})}
\item
  M is a flat A-module.
\item
  \(\operatorname{Tor}_1^{\mathbf{A}}(\mathbf{N},\mathbf{M}) = 0\) for every A-module \(\mathbf{N}\) annihilated by \(\Im\) .
\item
  M/3M is a flat (A/3)-module and the canonical mapping \(\Im \otimes_{\mathbf{A}} \mathbf{M} \to \Im \mathbf{M}\) is bijective (the latter condition being equivalent to \(\operatorname{Tor}_{1}^{\mathbf{A}}(\mathbf{A}/\Im, \mathbf{M}) = 0\) by virtue of the relation \(\operatorname{Tor}_{1}^{\mathbf{A}}(\mathbf{A}, \mathbf{M}) = 0\) and the exact sequence

\[
\operatorname{Tor} _ {1} ^ {\mathbf {A}} (\mathrm{A}, \mathrm{M}) \rightarrow \operatorname{Tor} _ {1} ^ {\mathbf {A}} (\mathrm{A} / \Im , \mathrm{M}) \rightarrow \Im \otimes_ {\mathrm{A}} \mathrm{M} \rightarrow \mathrm{M}).
\]

\item
  M/3M is a flat (A/3)-module and the canonical homomorphism

\[
\gamma_ {\mathbf {M}} \colon \operatorname{gr} (\mathbf {A}) \otimes_ {\mathbf {g r} _ {0} (\mathbf {A})} \operatorname{gr} _ {0} (\mathbf {M}) \rightarrow \operatorname{gr} (\mathbf {M})
\]

is bijective (property (GR) of § 2, no. 8).

\item
  For all \(n \geqslant 1\) , \(\mathbf{M} / \Im^n\mathbf{M}\) is a flat \((\mathbf{A} / \Im^n)\) -module.
\end{enumerate}

\[
(i) \Rightarrow (i i) \Leftrightarrow (i i i) \Rightarrow (i v) \Leftrightarrow (v).
\]

If further \(\mathfrak{I}\) is nilpotent or if A is Noetherian and M is ideally Hausdorff, properties (i), (ii), (iii), (iv) and (v) are equivalent.

Remark. If A/3 is a field (as often happens in applications) the condition ``M/3M is a flat (A/3)-module'' holds automatically for every A-module M, which simplifies the statement of properties (iii) and (iv); moreover, in this case, property (v) is equivalent to saying that M/3 \(^{n}\) M is a free (A/3 \(^{n}\) )-module for every integer n ≥ 1 (Chapter II, § 3, no. 2, Corollary 2 to Proposition 5).

\subsection{Proof of the flatness criterion}

\textbf{(A) The implications (i) \(\Rightarrow\) (ii) \(\Leftrightarrow\) (iii)}\par

The implication (i) \(\Rightarrow\) (ii) is immediate (Chapter I, § 4). The equivalence (ii) \(\Leftrightarrow\) (iii) is a special case of Chapter I, § 4, Proposition 2 applied to R = A, S = A/3, F = M, E = N, taking account of the fact that being given an (A/3)-module structure on N is equivalent to being given an A-module structure under which N is annihilated by 3.

Remark (1). Condition (ii) is also equivalent to the following:

(ii') \(\operatorname{Tor}_1^{\mathbf{A}}(\mathbf{N},\mathbf{M}) = 0\) for every A-module N annihilated by a power of \(\Im\) . Clearly (ii') implies (ii). Conversely, if (ii) holds, then in particular \(\operatorname{Tor}_1^{\mathbf{A}}(\Im^n\mathrm{N} / \Im^{n + 1}\mathrm{N},\mathbf{M}) = 0\) for all \(n\) ; from the exact sequence

\[
0 \rightarrow \Im^ {n + 1} N \rightarrow \Im^ {n} N \rightarrow \Im^ {n} N / \Im^ {n + 1} N \rightarrow 0
\]

we derive the exact sequence

\[
\operatorname{Tor} _ {1} ^ {\mathbf {A}} \left(\Im^ {n + 1} \mathrm{N}, \mathrm{M}\right)\rightarrow \operatorname{Tor} _ {1} ^ {\mathbf {A}} \left(\Im^ {n} \mathrm{N}, \mathrm{M}\right)\rightarrow \operatorname{Tor} _ {1} ^ {\mathbf {A}} \left(\Im^ {n} \mathrm{N} / \Im^ {n + 1} \mathrm{N}, \mathrm{M}\right)
\]

and, as there exists an integer \(m\) such that \(\mathfrak{S}^m\mathrm{N} = 0\) , we deduce by descending induction on \(n\) that \(\operatorname{Tor}_1^{\mathbf{A}}(\mathfrak{S}^n\mathrm{N},\mathbf{M}) = 0\) for all \(n\leqslant m\) and in particular for \(n = 0\) .

It follows from this that if \(\Im\) is nilpotent, (ii) implies (i), for (ii') then means that \(\operatorname{Tor}_1^{\mathbf{A}}(\mathbf{N},\mathbf{M}) = 0\) for every A-module \(\mathbf{N}\) and hence that \(\mathbf{M}\) is flat (Chapter I, §4).

\textbf{(B) Let us prove the following proposition:}\par

PROPOSITION 1. Let A be a commutative ring, \(\mathfrak{I}\) an ideal of A and M an A-module. The following conditions are equivalent:

\begin{enumerate}
\def\labelenumi{(\alph{enumi})}
\item
  For all \(n \geqslant 1\) , \(\operatorname{Tor}_1^{\mathbf{A}}(\mathbf{A}/\mathfrak{I}^n, \mathbf{M}) = 0\) .
\item
  For all \(n \geqslant 1\) , the canonical homomorphism
\end{enumerate}

\[
\theta_ {n} \colon \mathfrak {I} ^ {n} \otimes_ {\mathbf {A}} \mathbf {M} \rightarrow \mathfrak {I} ^ {n} \mathbf {M}
\]

is bijective.

Moreover these conditions imply:

\begin{enumerate}
\def\labelenumi{(\alph{enumi})}
\setcounter{enumi}{2}
\tightlist
\item
  The canonical homomorphism \(\gamma_{\mathrm{M}}: \operatorname{gr}(\mathbf{A}) \otimes_{\mathbf{g}_{0}(\mathbf{A})} \operatorname{gr}_{0}(\mathbf{M}) \to \operatorname{gr}(\mathbf{M})\) is bijective. Conversely, if \(\Im\) is nilpotent, (c) implies (a) and (b).
\end{enumerate}

The equivalence of (a) and (b) follows from the exact sequence

\[
0 = \operatorname{Tor} _ {1} ^ {\mathbf {A}} (\mathrm{A}, \mathrm{M}) \rightarrow \operatorname{Tor} _ {1} ^ {\mathbf {A}} (\mathrm{A} / \Im^ {n}, \mathrm{M}) \rightarrow \Im^ {n} \otimes_ {\mathbf {A}} \mathrm{M} \rightarrow \mathrm{M}.
\]

Consider next the diagram

\[
\begin{array}{c} \mathfrak {I} ^ {n + 1} \otimes_ {\mathrm{A}} \mathrm{M} \longrightarrow \mathfrak {I} ^ {n} \otimes_ {\mathrm{A}} \mathrm{M} \longrightarrow (\mathfrak {I} ^ {n} / \mathfrak {I} ^ {n + 1}) \otimes_ {\mathrm{A}} (\mathrm{M} / \mathfrak {I M}) \longrightarrow 0 \\ \theta_ {n + 1} \Bigg \downarrow \quad \theta_ {n} \Bigg \downarrow \quad \gamma_ {n} \Bigg \downarrow \\ 0 \longrightarrow \mathfrak {I} ^ {n + 1} \mathrm{M} \longrightarrow \mathfrak {I} ^ {n} \mathrm{M} \longrightarrow \operatorname{gr} _ {n} (\mathrm{M}) \longrightarrow 0 \end{array}\tag{1}
\]

where we note that \((\mathfrak{I}^n / \mathfrak{I}^{n+1}) \otimes_{\mathbf{A}} (\mathbf{M} / \mathfrak{I}\mathbf{M})\) is canonically identified with \((\mathfrak{I}^n / \mathfrak{I}^{n+1}) \otimes_{\mathbf{A} / \mathfrak{S}} (\mathbf{M} / \mathfrak{I}\mathbf{M})\) . This diagram is commutative by definition of \(\gamma_n\) and its rows are exact. If (b) holds, \(\theta_n\) and \(\theta_{n+1}\) are bijective and so therefore is \(\gamma_n\) by definition of cokernel, hence (b) implies (c). Conversely, assuming that \(\mathfrak{I}\) is nilpotent, let us show that (c) implies (b); we shall argue by descending induction on \(n\) , since \(\mathfrak{I}^n \otimes_{\mathbf{A}} \mathbf{M} = \mathfrak{I}^n\mathbf{M} = 0\) for \(n\) sufficiently large. Suppose then that in diagram (1), \(\gamma_n\) and \(\theta_{n+1}\) are bijective; then so is \(\theta_n\) by virtue of Chapter I, § 1, no. 4, Corollary 1 to Proposition 2.

\textbf{(C) The implication (ii) \(\Rightarrow\) (iv)}\par

If (ii) holds, so does (ii') by Remark 1; Proposition 1 then shows that \(\gamma_{\mathbf{M}}\) is an isomorphism. On the other hand, we already know that (ii) implies (iii) and hence M/Im is a flat (A/Im)-module, which completes the proof that (ii) implies (iv).

Remark (2). Proposition 1 shows that, if \(\mathfrak{S}\) is nilpotent, (iv) implies (iii); taking account of Remark 1, we have therefore proved in this case that (i), (ii), (iii) and (iv) are equivalent.

\textbf{(D) The equivalence (iv) \(\Leftrightarrow\) (v)}\par

For all \(n \geqslant 1\) , M/ \(\Im^{n}M\) has a canonical (A/ \(\Im^{n}\) )-module structure. If it is filtered by the \((\Im/\Im^{n})\) -adic filtration, it is immediate that \(\mathrm{gr}_{m}(\mathrm{M}/\Im^{n}\mathrm{M}) = \mathrm{gr}_{m}(\mathrm{M})\) if m \textless{} n and \(\mathrm{gr}_{m}(\mathrm{M}/\Im^{n}\mathrm{M}) = 0\) if \(m \geqslant n\) . For all \(k \geqslant 1\) , let \(A_{k} = A/\Im^{k}\) , \(\Im_{k} = \Im/\Im^{k} M_{k} = M/\Im^{k} M\) ; let \((\mathrm{iv})_{k}\) (resp. \((\mathrm{v})_{k}\) ) denote the assertion derived from (iv) (resp. (v)) by replacing A, \(\Im\) , M by \(A_{k}\) , \(\Im_{k}\) , \(M_{k}\) . It follows from what has just been said that (iv) is equivalent to ``for all \(k \geqslant 1\) , \((\mathrm{iv})_{k}\) '' and obviously (v) is equivalent to ``for all \(k \geqslant 1\) , \((\mathrm{v})_{k}\) ''. Then it will suffice to establish the equivalence \((\mathrm{iv})_{k} \Leftrightarrow (\mathrm{v})_{k}\) for all k or also to show that (iv) \(\Leftrightarrow (\mathrm{v})\) when \(\Im\) is nilpotent. Now (Remark 2) we have seen that in that case (iv) is equivalent to (i). As M/ \(\Im^{n}M\) is isomorphic to M \(\otimes_{A} (A/\Im^{n})\) , (i) implies (v) (Chapter I, §2, no. 7, Corollary 2 to Proposition 8); moreover clearly (v) then implies (i). We have therefore shown the equivalence (iv) \(\Leftrightarrow (\mathrm{v})\) in all cases and also that of all the properties of the theorem in the case where \(\Im\) is nilpotent.

\textbf{(E) The implication \((\mathbf{v})\Rightarrow (\mathbf{i})\) when \(\mathbf{A}\) is Noetherian and \(\mathbf{M}\) ideally Hausdorff}\par

It is sufficient to prove that for every ideal \(\mathfrak{a}\) of \(\mathbf{A}\) the canonical mapping \(j\colon \mathfrak{a} \otimes_{\mathbf{A}} \mathbf{M} \to \mathbf{M}\) is injective (Chapter I, § 2, no. 3, Proposition 1). Let \(x \in \operatorname{Ker} j\) ; as \(\mathfrak{a} \otimes_{\mathbf{A}} \mathbf{M}\) is Hausdorff with the \(\mathfrak{S}\) -adic topology, it suffices to verify that, for every integer \(n > 0\) , \(x \in \mathfrak{S}^n(\mathfrak{a} \otimes_{\mathbf{A}} \mathbf{M})\) . Let \(f\colon \mathfrak{S}^n\mathfrak{a} \to \mathfrak{a}\) be the canonical injection; it suffices to show that \(x \in \operatorname{Im}(f \otimes 1_{\mathbf{M}})\) ; for if \(b \in \mathfrak{S}^n\) , \(a \in \mathfrak{a}\) and \(m \in \mathbf{M}\) , the image under \(f \otimes 1_{\mathbf{M}}\) of the element \((ba) \otimes m\) of \((\mathfrak{S}^n\mathfrak{a}) \otimes_{\mathbf{A}} \mathbf{M}\) is the element \((ba) \otimes m = b(a \otimes m)\) of \(\mathfrak{a} \otimes_{\mathbf{A}} \mathbf{M}\) and hence \(\operatorname{Im}(f \otimes 1_{\mathbf{M}}) \subset \mathfrak{S}^n(\mathfrak{a} \otimes_{\mathbf{A}} \mathbf{M})\) . By virtue of Krull's Theorem (§ 3, no. 2, Theorem 2), there exists an integer \(k\) such that \(\mathfrak{a}_k = \mathfrak{a} \cap \mathfrak{S}^k \subset \mathfrak{S}^n\mathfrak{a}\) ; if \(i: \mathfrak{a}_k \to \mathfrak{a}\) is the canonical injection, it will then be sufficient to show that \(x \in \operatorname{Im}(i \otimes 1_M)\) . Now, denoting by \(p: \mathfrak{a} \to \mathfrak{a}/\mathfrak{a}_k\) and \(h: \mathfrak{a}/\mathfrak{a}_k \to \mathrm{A}/\mathfrak{S}^k\) the canonical mappings, there is a commutative diagram

\[
\begin{array}{c} \mathfrak {a} _ {k} \otimes_ {\mathrm{A}} \mathrm{M} \xrightarrow {i \otimes 1 _ {\mathrm{M}}} \mathfrak {a} \otimes_ {\mathrm{A}} \mathrm{M} \xrightarrow {p \otimes 1 _ {\mathrm{M}}} (\mathfrak {a} / \mathfrak {a} _ {k}) \otimes_ {\mathrm{A}} \mathrm{M} \longrightarrow 0 \\ j \Biggl \downarrow \\ \mathrm{M} \quad \longrightarrow \quad (\mathrm{A} / \mathfrak {I} ^ {k}) \otimes_ {\mathrm{A}} \mathrm{M} \end{array}
\]

in which the first row is exact. It suffices to prove that \(x \in \operatorname{Ker}(p \otimes 1_{\mathbf{M}})\) and, as \(x \in \operatorname{Ker} j\) by hypothesis, it will suffice to verify that the mapping \(h \otimes 1_{\mathbf{M}}\) is injective. Now, it may also be written (Algebra, Chapter II, § 3, no. 6, Corollary 3 to Proposition 6)

\[
h \otimes 1 _ {\mathbf {M} / \Im^ {k} \mathbf {M}}: (a / a _ {k}) \otimes_ {\mathbf {A} / \Im^ {k}} (\mathbf {M} / \Im^ {k} \mathbf {M}) \rightarrow \mathbf {M} / \Im^ {k} \mathbf {M}
\]

and, as \(h\) is injective and, by (v), \(\mathbf{M} / \Im^k\mathbf{M}\) is a flat \((\mathbf{A} / \Im^k)\) -module, this completes the proof.

\subsection{Applications}

PROPOSITION 2. Let A be a commutative ring, \(\mathfrak{S}\) an ideal of A and B a commutative Noetherian A-algebra such that B is contained in the Jacobson radical of B. Then every finitely generated B-module M is an ideally Hausdorff A-module with respect to \(\mathfrak{S}\) .

We shall see more generally that for every finitely generated A-module N, \(N \otimes_{A} M\) is Hausdorff with the \(\Im\) -adic topology. For \(\mathrm{N}_{(\mathrm{B})} = \mathrm{N} \otimes_{\mathrm{A}} \mathrm{B}\) is a finitely generated B-module and the B-module \(N \otimes_{A} M\) is canonically identified with \(\mathrm{N}_{(\mathrm{B})} \otimes_{\mathrm{B}} M\) by virtue of the associativity of the tensor product. Let L be the Jacobson radical of B; as \(\Im B\) is contained in L, the \(\Im\) -adic topology on \(N \otimes_{A} M\) is therefore identified with a finer topology than the \(\Im\) -adic topology on \(\mathrm{N}_{(\mathrm{B})} \otimes_{\mathrm{B}} M\) ; but this latter topology is Hausdorff since \(\mathrm{N}_{(\mathrm{B})} \otimes_{\mathrm{B}} M\) is a finitely generated B-module (no. 1, Example 1), whence the conclusion.

PROPOSITION 3. Let A be a commutative ring, B a commutative A-algebra, \(\mathfrak{I}\) an ideal of A and M a B-module. Suppose that B is a Noetherian ring and a flat A-module and that M is ideally Hausdorff with respect to \(\mathfrak{S}\) B. The following conditions are equivalent:

\begin{enumerate}
\def\labelenumi{(\alph{enumi})}
\item
  M is a flat B-module.
\item
  \(\mathbf{M}\) is a flat A-module and \(\mathbf{M} / \Im \mathbf{M} = \mathbf{M} / (\Im \mathbf{B})\mathbf{M}\) is a flat (B/3B)-module.
\end{enumerate}

If further the canonical homomorphism \(\mathbf{A} / \Im \to \mathbf{B} / \Im \mathbf{B}\) is bijective, conditions (a) and (b) are also equivalent to:

\begin{enumerate}
\def\labelenumi{(\alph{enumi})}
\setcounter{enumi}{2}
\tightlist
\item
  M is a flat A-nodule.
\end{enumerate}

Condition (a) implies (b) by Chapter I, §2, no. 7, Corollaries 2 and 3 to Proposition 8 and the fact that M/Im is isomorphic to M ⊗B (B/ImB). Suppose condition (b) holds; to show that M is a flat B-module, we shall apply Theorem 1 of no. 2 with A replaced by B and S by SIB. It will therefore be sufficient to show that the canonical mapping \(f\colon \mathfrak{B} \otimes_{\mathbf{B}} \mathbf{M} \to \mathfrak{M}\) is injective. Let \(f_1\) be the canonical mapping \(\mathfrak{I} \otimes_{\mathbf{A}} \mathbf{B} \to \mathfrak{B}\) and \(f_2\) the canonical isomorphism \(\mathfrak{I} \otimes_{\mathbf{A}} \mathbf{M} \to (\mathfrak{I} \otimes_{\mathbf{A}} \mathbf{B}) \otimes_{\mathbf{B}} \mathbf{M}\) ; \(f \circ (f_1 \circ 1_{\mathbf{M}}) \circ f_2\) is the canonical mapping \(f'\colon \mathfrak{I} \otimes_{\mathbf{A}} \mathbf{M} \to \mathfrak{IM}\) , as is easily verified. Now \(f'\) is an isomorphism since M is a flat A-module, whilst \(f_1\) is an isomorphism because B is flat over A; \(f\) is then an isomorphism.

Let \(\rho: \mathring{\mathrm{A}}/\mathfrak{I} \to \mathrm{B}/\mathfrak{I}\mathrm{B}\) be the canonical homomorphism; the (A/ \(\mathfrak{I}\) )-module structure on M/ \(\mathfrak{I}\) M derived by means of \(\rho\) from its (B/ \(\mathfrak{I}\) B)-module structure is isomorphic to that on M \(\otimes_{\mathrm{A}} (\mathrm{A}/\mathfrak{I})\) . Then it follows that, if M is a flat A-module, M/ \(\mathfrak{I}\) M is a flat (A/ \(\mathfrak{I}\) )-module and hence also a flat (B/ \(\mathfrak{I}\) B)-module if \(\rho\) is an isomorphism; we have thus proved that (c) \(\Rightarrow\) (b) in that case.

COROLLARY. Let A be a commutative Noetherian ring, \(\Im\) an ideal of A, \(\hat{\mathbf{A}}\) the Hausdorff completion of A with respect to the \(\Im\) -adic topology and M an ideally Hausdorff \(\hat{\mathbf{A}}\) -module with respect to \(\Im\hat{\mathbf{A}}\) . For M to be a flat A-module, it is necessary and sufficient that M be a flat \(\hat{\mathbf{A}}\) -module.

We know in fact that \(\hat{\mathbf{A}}\) is a Noetherian ring ( \(\S 3\) , no. 4, Proposition 8) and a flat A-module ( \(\S 3\) , no. 4, Theorem 3), that \(\Im \hat{\mathbf{A}} = \hat{\mathfrak{I}}\) ( \(\S 2\) , no. 12, Proposition 16) and that the canonical homomorphism \(\mathrm{A} / \Im \to \hat{\mathrm{A}} / \hat{\mathfrak{I}}\) is bijective ( \(\S 2\) , no. 12, Proposition 15); Proposition 3 can therefore be applied.

PROPOSITION 4. Let A and B be two commutative Noetherian rings, h: A → B a ring homomorphism, ∑ an ideal of A and ∑ an ideal of B containing ∑B and contained in the Jacobson radical of B. Let A be the Hausdorff completion of A with respect to the ∑-adic topology and B the Hausdorff completion of B with respect to the ∑-adic topology; h is continuous with these topologies and h: A → B therefore makes B into an A-algebra. Let M be a finitely generated B-module and M its Hausdorff completion with respect to the ∑-adic topology; the following properties are equivalent:

\begin{enumerate}
\def\labelenumi{(\alph{enumi})}
\item
  M is a flat A-module.
\item
  \(\hat{\mathbf{M}}\) is a flat A-module.
\item
  \(\hat{\mathbf{M}}\) is a flat \(\hat{\mathbf{A}}\) -module.
\end{enumerate}

As B with the \(\mathfrak{L}\) -adic topology is a Zariski ring, \(\hat{\mathbf{B}}\) is a faithfully flat B-module (§ 3, no. 5, Proposition 9) and \(\hat{\mathbf{M}}\) is canonically isomorphic to \(\mathbf{M} \otimes_{\mathbf{B}} \hat{\mathbf{B}}\) (§ 3, no. 4, Theorem 3); it is immediately verified that this canonical isomorphism is an isomorphism of the A-module structure on \(\hat{\mathbf{M}}\) onto the A-module structure on \(\mathbf{M} \otimes_{\mathbf{B}} \hat{\mathbf{B}}\) derived from that on M. Applying Proposition 4 of Chapter I, § 3, no. 2 with R replaced by B, S by A, E by \(\hat{\mathbf{B}}\) , F by M, we see that for M to be a flat A-module, it is necessary and sufficient that \(\hat{\mathbf{M}}\) be a flat A-module. Moreover, \(\hat{\mathbf{M}}\) is a finitely generated \(\hat{\mathbf{B}}\) -module and \(\mathfrak{J}\hat{\mathbf{B}}\) is contained in \(\hat{\mathfrak{L}} = \mathfrak{L}\hat{\mathbf{B}}\) and hence in the Jacobson radical of \(\hat{\mathbf{B}}\) (§ 3, no. 4, Proposition 8); therefore \(\hat{\mathbf{M}}\) is an ideally Hausdorff \(\hat{\mathbf{A}}\) -module with respect to \(\mathfrak{J}\hat{\mathbf{A}}\) (Proposition 2). Conditions (b) and (c) are therefore equivalent by the Corollary to Proposition 3.


\exercisesection{1}

\begin{enumerate}
\def\labelenumi{\arabic{enumi}.}
\tightlist
\item
  Let A be a graded commutative ring of type Z and \((A_{n})\) its graduation; we set \(A^{\geq} = \bigoplus_{n \geq 0} A_{n}, A^{\leq} = \bigoplus_{n \leqslant 0} A_{n}\) ; these are graded subrings of A.

\begin{enumerate}
\def\labelenumii{(\alph{enumii})}
\item
  For every homogeneous element \(f\) of \(\mathbf{A}\) of degree \(d\) , the ring of fractions \(\mathbf{A}_f\) (corresponding to the multiplicative subset consisting of the \(f^n\) , where \(n \geqslant 0\) ) has a canonical graded ring structure (Algebra, Chapter II, §11); we denote by \(\mathbf{A}_{(f)}\) the subring of \(\mathbf{A}_f\) consisting of elements of degree 0. Show that, if \(d > 0\) , then \((\mathbf{A}^{\geq})_f = \mathbf{A}_f, \mathbf{A}_{(f)}\) is isomorphic to \(\mathbf{A}^{(d)} / (f - 1)\mathbf{A}^{(d)}\) and \(((\mathbf{A}_f)^{\geq})_{f/1}\) is a graded ring isomorphic to \(\mathbf{A}_f\) .
\item
  Let the polynomial ring \(\mathbf{B} = \mathbf{A}[\mathbf{X}] = \mathbf{A} \otimes_{\mathbf{Z}} \mathbf{Z}[\mathbf{X}]\) be given the graduation the tensor product of those on \(\mathbf{A}\) and \(\mathbf{Z}[\mathbf{X}]\) , a graduation compatible with the ring structure on \(\mathbf{B}\) . Show that, if \(d > 0\) , \(\mathrm{B}_{(f)}\) is isomorphic to \((\mathrm{A}_f)^{\leqslant}\) and \((\mathrm{A}^{(d)})_f\) is a graded ring isomorphic to \(\mathrm{A}_{(f)} \otimes_{\mathbf{Z}} \mathbf{Z}[\mathbf{X}, \mathbf{X}^{-1}]\) .
\item
  Let \(g\) be another homogeneous element of \(A\) of degree \(e\) . Show that, if \(d > 0\) and \(e > 0\) , \(A_{(fg)}\) is isomorphic to \((A_{(f)})_{g^d / f^e}\) .\\
\item[* (d)] Suppose \(d > 0\) and \(A_n = \{0\}\) for \(n < 0\) . Show that, if \(A\) is Noetherian, so is \(A_{(f)}\) (use §2, no. 10, Corollary 4 to Theorem 2). *
\end{enumerate}

\item
  Let A be a graded commutative ring of type Z such that \(A_{n}=0\) for n\textless0. For all \(n\geqslant0\) , let \(A_{[n]}\) denote the graded ideal \(\bigoplus_{m\geqslant n}A_{m}\) of A; the A-algebra \(A^{\sharp}=\bigoplus_{n\geqslant0}A_{[n]}\) is a ring graded by the \(A_{[n]}=A_{n}^{\sharp}\) ; for all \(f\in A_{d}(d>0)\) , let \(f^{\sharp}\) denote the element of \(A^{\sharp}\) whose component are zero except for that of degree d, which is equal to f.

\begin{enumerate}
\def\labelenumii{(\alph{enumii})}
\item
  Show that, if \(f \in \mathbf{A}_d\) , where \(d > 0\) , the ring \(\mathrm{A}_{(f^{\sharp})}^{\sharp}\) is isomorphic to \((\mathbf{A}_f)^{\geqslant}\) (notation of Exercise 1).
\item
  For \(\mathbf{A}^{\natural}\) to be a finitely generated A-algebra, it is necessary and sufficient that \(\mathbf{A}\) be a finitely generated \(\mathbf{A}_0\) -algebra.
\item
  In order that \(\mathbf{A}_{n+1}^{\natural} = \mathbf{A}_1^{\natural}\mathbf{A}_n^{\natural}\) (resp. \(\mathbf{A}_n^{\natural} = (\mathbf{A}_1^{\natural})^n\) ) for \(n \geqslant n_0\) , it is necessary and sufficient that \(\mathbf{A}_{n+1} = \mathbf{A}_1\mathbf{A}_n\) (resp. \(\mathbf{A}_n = (\mathbf{A}_1)^n\) ) for \(n \geqslant n_0\) .
\end{enumerate}

\item
  Let K be a field, B the polynomial ring \(\mathrm{K}[X,Y]\) with the graduation defined by the total degree, C the graded subring of B generated by X and \(Y^{2}\) and a the graded ideal of C generated by \(Y^{4}\) . Show that, in the graded ring A = C/a, \(A_{n+1} = A_{1}A_{n}\) for \(n \geqslant 2\) but \(A_{n} \neq (A_{1})^{n}\) for \(n \geqslant 6\) .
\end{enumerate}
